
\section{Introduction}

The rapid advancement of large-scale code generation models (LCGMs) has enabled them to address a broad range of general programming tasks~\cite{chen2021evaluating, openai2023gpt4, team2023gemini, touvron2023llama2}. Although these advances are promising, they also raise critical concerns regarding performance evaluation. A benchmark that enables accurate and objective assessment is indispensable for both users and developers. Consequently, numerous benchmarks have been introduced~\cite{chen2021evaluating, jain2024livecodebench, zhuo2024bigcodebench}, most of which rely on static, ground-truth-based evaluation.

\subsection{Data Contamination}

``When a measure becomes a target, it ceases to be a good measure.'' Achieving trustworthy evaluation is challenging. One of the most significant obstacles is \emph{data contamination}, which arises when test data is included in a model's training set, leading to an overestimation of its true performance. Because of this inflation, the observed gains do not generalize to real-world programming tasks, the benchmark's utility declines, and developers lose confidence in the reported numbers. In effect, every static benchmark has a finite useful lifetime.

Manually refreshing static benchmarks is an insufficient remedy for data contamination. The problem is rooted in the static and open-source nature of contemporary benchmarks, which allows evaluation tasks to be readily incorporated into training corpora, thereby artificially inflating accuracy metrics~\cite{yang2023rethinking}. Rebuilding a benchmark from scratch is straightforward in principle, but problematic in practice. First, designing and disseminating a new benchmark---especially one that targets a specialized domain---requires substantial time and human effort. Second, although frequent benchmark updates seek to mitigate this risk, the pace of contamination often exceeds the multi-month cycles required for formal benchmark releases. Consequently, the rapid fine-tuning of models---often completed within days---can render newly released datasets such as KernelBench~\cite{lange2025robustagenticcudakernel} or historical problems in LiveCodeBench~\cite{jain2024livecodebench} obsolete almost immediately. Therefore, manual updates to static benchmarks cannot adequately address data contamination.
To tackle this issue, existing research has explored two primary directions. The first direction is \emph{task perturbation}, which translates, rewrites, adapts, or merges existing benchmark tasks so that the updated tasks differ from the data seen during training. The second direction is \emph{contamination detection}, which uses predefined features to detect whether a model has been contaminated by a benchmark; some methods even leverage detection results to mitigate the contamination.

\subsection{Limitations of Existing Methods}

Both directions exhibit fundamental limitations. First, traditional task-generation and perturbation methods often fail to produce sufficiently challenging tasks. Approaches such as RECODE~\cite{wang-etal-2023-recode} primarily perturb task descriptions while leaving canonical solutions unchanged and therefore retrievable from model parameters. Although PPM~\cite{10.1145/3643780} transforms code outputs, these variants are often too simplistic for modern LCGMs. More importantly, current generation methods---such as sequentially combining independent problems into a single task---are easily decomposed and solved by advanced models. Ultimately, creating complex, diverse, and verified task variants still requires substantial manual effort to ensure the correctness of test cases, thereby lacking the scalability necessary for rapid iteration of LLMs.

Second, methods based on contamination detection are constrained by inaccurate leakage patterns and limited applicability. Many effective detection techniques~\cite{nguyen2025rn} require access to internal model parameters or log-probabilities, which renders them inapplicable to widely used closed-source LLMs such as GPT, Claude, and Gemini. Even in output-based detection, methods that measure similarity across multiple outputs~\cite{dong-etal-2024-generalization} frequently yield false positives. The fundamental flaw in this reasoning is that, for simple but fundamental problems, nearly all correct solutions are inherently similar; incorrectly flagging these as leaked results in a biased and ineffective evaluation of a model's true capability. More importantly, most contamination detection methods cannot be directly applied to reliability assessment because they cannot determine how the model actually performs on test sets that have not been leaked.

\subsection{Our Approach}

To address these limitations, we introduce \tool{}, an automated pipeline for adapting programming tasks. Inspired by the ability of models to solve concrete problems, we hypothesize that they can also adapt problems with high quality. We exploit this capability to design a fully automated pipeline. First, we provide a model with the canonical solution of a programming task and pose an open-ended refactoring request, asking the model to modify the code's functionality. Next, we construct a code-generation prompt for the mutated code, including the problem description and test cases, thereby forming a new, self-contained programming task. Finally, the pipeline equips the task with the corresponding test harness. In the following sections, we discuss how our method mitigates data contamination and empirically demonstrate its effectiveness.

\myparagraph{A High-Level Approach to Our Evaluation.} The objective of our method is to achieve automated evaluation while mitigating contamination. We follow the idea of task- generation based methods, but realize it by leveraging LLMs. First, we formulate an open-ended question that prompts models to mutate code selected from the current benchmark in order to transform its functionality. Then, we construct a novel programming task, which primarily comprises the prompt that guides the model to generate the target code and the test code based on the mutated code, and we use it to evaluate the current model. We repeat the above procedure several times and finally calculate each model's score from its answers and the quality of the mutated code.


\hzbc{Put challenge here better or 2.1???    Here I will supply the criteria the quality of the mutated task }
It is clear that the main challenge lies in the first step---how can we ensure that the mutated task is able to mitigate data contamination? One straightforward approach would be to ask the model to generate a program together with the problem statement and test cases directly. However, this would raise several issues: the generated test cases may be inconsistent with the code, the problem description may be ambiguous, and the resulting tasks may be either too easy or computationally infeasible. Instead, we pursue an alternative route: we mutate an existing verified solution and derive the new task and test cases directly from the mutated code. This guarantees consistency and tractability by construction.

\label{sec:intro}
